% UNR accessible thesis / dissertation — TeXLib template.
% Compile with lualatex. \DocumentMetadata MUST be the first line (it is what
% switches on PDF tagging, and it can only do so before \documentclass).
\DocumentMetadata{
	lang          = en,
	tagging       = on,
	tagging-setup = { math/setup={mathml-SE}, table/header-rows=1 },
	pdfstandard   = { ua-2, a-4f },
}
\documentclass{thesis}

% --- Degree / committee metadata ---------------------------------------------
\title{An Accessible Approach to Something Important}
\author{Ada N. Student}
\doctype{thesis}                 % thesis (M.S.) or dissertation (Ph.D.)
\gradprogram{Mathematics}
\gradadvisor{Pat D. Advisor, Ph.D.}
\graddate{May, 2027}
\committeemember{Pat D. Advisor, Ph.D.}{Advisor}
\committeemember{Sam K. Reader, Ph.D.}{Committee Member}
\committeemember{Lee Q. Faculty, Ph.D.}{Graduate School Representative}

\addbibresource{thesis-refs.bib}

\begin{document}

% --- Front matter ------------------------------------------------------------
\makeUNRtitlepage
\makecopyrightpage
\committeeapprovalpage

\frontmatter

\begin{frontmatterpage}{Abstract}
This thesis develops an accessible approach to a problem of genuine importance,
and demonstrates that the resulting document is both mathematically rigorous and
readable by assistive technology. We show that the operator $T$ satisfies
$T(x) = \ThesisMathIntent{inverse-sine}{\sin^{-1}}(x)$ for all admissible $x$.
\end{frontmatterpage}

\begin{frontmatterpage}{Acknowledgements}
I thank my advisor, my committee, and the veraPDF validator.
\end{frontmatterpage}

\tableofcontents
\listoffigures

\mainmatter

% --- Body --------------------------------------------------------------------
\chapter{Introduction}
This is the introduction. Cross-references work: see \cref{thm:main} below.

\begin{definition}[Admissible]
A point $x$ is \emph{admissible} if $x > 0$ and $x \neq 1$.
\end{definition}

\begin{theorem}[Main result]\label{thm:main}
For every admissible $x$, the identity $\int_0^x t \, dt = \tfrac{1}{2}x^2$ holds.
\end{theorem}

\begin{proof}
Integrate directly.
\end{proof}

\chapter{Methods}
The methods chapter continues the argument. \Cref{thm:main} is the key tool,
following the classic treatment~\autocite{rudin1976}.

\begin{example}
Take $x = 2$: then $\int_0^2 t\,dt = 2$, and indeed $\tfrac{1}{2}(2)^2 = 2$.
\end{example}

\makereferences

\end{document}
